% !TeX root = ../main.tex
\section{Расчет каскада на нелинейной модели методом малого и большого сигнала}

\subsection{Малосигнальный частотный анализ каскада с цепями смещения}

Исследование линейных параметров каскада на основе нелинейной модели SPICE (\texttt{NE68133Package}) позволяет учесть совместное влияние паразитных параметров корпусирования, нелинейных переходов и реальной схемы питания по постоянному току. Принципиальная схема усилителя с цепями термостабилизированного резистивного смещения и согласующими звеньями приведена на рис.~\ref{fig:nl_schematic}.

\begin{figure}[H]
    \centering
    \includegraphics[width=0.82\textwidth]{03_sch.png}
    \caption{Схема усилителя на нелинейной модели NE68133 с питанием от единого источника DC}
    \label{fig:nl_schematic}
\end{figure}

Подача напряжения смещения на базу транзистора реализована посредством резистивного делителя $R_{35} = 8{,}25\text{ кОм}$ и $R_{36} = 5{,}1\text{ кОм}$, подключенного к основному источнику питания $V_7 = 2{,}5\text{ В}$. Разделительные емкости $C_{28} = 1{,}09\text{ пФ}$ и $C_{31} = 100\text{ пФ}$ исключают шунтирование постоянного потенциала базы через заземленную индуктивность входного согласования $L_{27} = 8{,}28\text{ нГн}$ и входной 50-омный тракт. Питание коллектора осуществляется через радиочастотный дроссель $L_{49} = 10\text{ нГн}$, а нагрузка развязана конденсатором $C_{51} = 10\text{ пФ}$.

В среде Ansys Designer был проведен частотный анализ (\textit{Linear Network Analysis}) в диапазоне $1{,}0\dots 6{,}0\text{ ГГц}$ с вычислением коэффициента шума (\textit{Enable Noise Calculation}). Результаты совместного моделирования $S$-параметров и коэффициента шума $NF$ представлены на рис.~\ref{fig:nl_s_nf}.

\begin{figure}[H]
    \centering
    \includegraphics[width=0.78\textwidth]{03_plot.png}
    \caption{Частотные зависимости $S$-параметров и коэффициента шума каскада на нелинейной модели}
    \label{fig:nl_s_nf}
\end{figure}

На центральной рабочей частоте $f_0 = 4{,}5\text{ ГГц}$ получены следующие характеристики:
\begin{itemize}
    \item Коэффициент согласования по выходу $\mathrm{dB}(S_{22}) \le -13\text{ дБ}$ с ярко выраженным резонансным минимумом $-18\text{ дБ}$ на частоте $5{,}1\text{ ГГц}$;
    \item Входное согласование $\mathrm{dB}(S_{11}) \approx -10\text{ дБ}$;
    \item Коэффициент шума каскада $NF \approx 3{,}5\text{ дБ}$;
    \item Коэффициент передачи $\mathrm{dB}(S_{21}) \approx 0\dots +0{,}85\text{ дБ}$, что объясняется дополнительным шунтирующим действием высокоомного базового делителя и технологической эмиттерной индуктивности обратной связи.
\end{itemize}

\subsection{Анализ компрессии коэффициента передачи методом гармонического баланса}

Для нахождения динамического диапазона усилителя и определения верхней границы линейности был проведен расчет методом однотонового гармонического баланса (\textit{Harmonic Balance 1-Tone}, сетап \texttt{PSweep}) на рабочей частоте $f_0 = 4{,}5\text{ ГГц}$. Мощность входного сигнала варьировалась в диапазоне от $-40$ до $+10\text{ дБм}$ с шагом $1\text{ дБм}$.

Зависимости выходной мощности на первой гармонике $\mathrm{dBm}(P_{\text{вых}})$ и коэффициента передачи $\mathrm{dB}(TG_{21})$ от уровня входной мощности приведены на рис.~\ref{fig:nl_p1db}.

\begin{figure}[H]
    \centering
    \includegraphics[width=0.78\textwidth]{03_plot2.png}
    \caption{Зависимость выходной мощности и коэффициента передачи от входного уровня (анализ сжатия усиления $P_{1\mathrm{dB}}$)}
    \label{fig:nl_p1db}
\end{figure}

По данным маркеров определены ключевые нелинейные параметры:
\begin{itemize}
    \item В области малого сигнала ($P_{\text{вх}} \le -30\text{ дБм}$) каскад работает строго в линейном режиме с коэффициентом передачи $G_0 \approx +0{,}85\text{ дБ}$;
    \item При увеличении мощности до $P_{\text{вх}} = -8{,}0\text{ дБм}$ наблюдается спад коэффициента передачи на $0{,}5\text{ дБ}$ ($\Delta G = -0{,}4986\text{ дБ}$, маркеры $m_1, m_2$);
    \item Точка компрессии коэффициента передачи на 1 дБ ($P_{1\mathrm{dB},\text{вх}}$) наступает при уровне входного воздействия $P_{\text{вх}} \approx -5{,}5\dots -6{,}0\text{ дБм}$, при этом уровень выходной мощности составляет $P_{1\mathrm{dB},\text{вых}} \approx -6{,}0\text{ дБм}$.
\end{itemize}

Таким образом, верхняя граница линейного динамического диапазона усилителя по входу составляет около $-6\text{ дБм}$, что гарантирует отсутствие искажений при типовых уровнях сигналов на входе малошумящего тракта.
\subsection{Расчет переходного процесса во временной области (Transient Analysis)}

Для оценки динамического отклика каскада во временной области при гармоническом воздействии был выполнен расчет переходных характеристик с использованием решателя \textit{Nexxim Transient Analysis}. В соответствии с заданием режим моделирования задавался следующими параметрами:
\begin{itemize}
    \item Частота входного синусоидального воздействия: $f = 1{,}0\text{ ГГц}$ (период колебаний $T = 1{,}0\text{ нс}$);
    \item Напряжение источника питания каскада: $U_{\text{пит}} = 2{,}0\text{ В}$;
    \item Амплитуда входного сигнала: $A_{\mathrm{tr}} = 0{,}01\text{ В}$ ($10\text{ мВ}$).
\end{itemize}

Интервал интегрирования был выбран равным $t_{\mathrm{stop}} = 10\text{ нс}$ с шагом дискретизации $\Delta t = 0{,}01\text{ нс}$ для обеспечения высокой точности аппроксимации формы волны. Полученные временные зависимости напряжений в узлах входа $V(\text{Input})$ и выхода $V(\text{Output})$ представлены на рис.~\ref{fig:transient_plot}.

\begin{figure}[H]
    \centering
    \includegraphics[width=0.82\textwidth]{04_transient.png}
    \caption{Осциллограммы входного и выходного напряжений каскада в режиме переходного процесса}
    \label{fig:transient_plot}
\end{figure}

По результатам анализа переходного процесса зафиксировано:
\begin{itemize}
    \item Время установления стационарного гармонического режима составляет менее $1{,}5\text{ нс}$, после чего форма выходного сигнала приобретает стабильный периодический характер без признаков самовозбуждения или низкочастотной релаксации;
    \item Постоянная составляющая напряжений стабилизирована на уровне напряжения питания $2{,}0\text{ В}$;
    \item Усилительный каскад сохраняет устойчивость при подаче динамического сигнала с заданной амплитудой.
\end{itemize}

